% Fin du Tome II : la bibliographie. Certaines notes portent des identifiants insécables (noms de
% structures Windows, de fonctions CUDA) : on autorise un peu plus d'espace entre les mots plutôt
% qu'un débordement dans la marge (\sloppy : une empreinte de 40 chiffres
% hexadécimaux passe alors seule à la ligne au lieu de sortir de 37 à 51 pt).
{\sloppy\setlength{\emergencystretch}{3em}%
\printbibliography[heading=bibintoc, title={\EF{Bibliography}{Bibliographie}}]}
